\chapter{Bridge}
\section{Purpose}

This case is not properly saying a validation test case but more an example to 
demonstrate a possible solution for user to manage bridge in TELEMAC-2D.
In this case, the choice was made to represent the bridge as an obstacle to 
the flow (like a dike in the river bed) where TELEMAC-2D will compute the 
overflowing discharge. Nevertheless, as the part of the discharge passing under 
the bridge is not neglictible and to also to allow the computation of the full 
hydrograph, the flow which pass under the bridge is take in account by an 
internal singularity, under the form of a couple of source and sink nodes.

\section{Geometry and mesh}

A river reach with it's major bed of 1000 m long crossed by a road in the middle. 
At the location of the main channel, it exists under the road 2 rectangular tubes
of 2,5 m width and 1,5 m height.

The mesh of the main channel and the road is regular (10 x 10 m).
The rest of the mesh is computed without any constraint by Blue Kenue

\begin{itemize}
\item  5000 triangular elements
\item  2626 nodes
\item  Maximum size range: 16.36 meters
\end{itemize}

The Figure~\ref{fig:bridge:mesh} shows the mesh of the study.
\begin{figure}
\centering
\includegraphicsmaybe{[width=.8\textwidth]}{../img/Mesh.png}
\caption{Mesh of the study}\label{fig:bridge:mesh}
\end{figure}

\section{Boundaries}

Lateral:

\begin{itemize}
\item  Solid walls with slip condition on both lateral parts of the major bed
\item  Imposed discharge upstream (left). The discharge vary in time from 0 at the beginning
 it increase to 20 m${}^{3}$/s and stay at this value between 100 s and 6100 s. Then the
 discharge increase to 120 m${}^{3}$/s and stay at this value between 6600 s and 17600 s. 
 After that the discharge decrease to 20 m${}^{3}$/s and stay at this level for the 
 rest of the simulation
\item  Imposed water level downstream (right). The value of 3 m is imposed
\end{itemize}

Bottom:

\begin{itemize}
\item  Strickler formula with friction coefficient = 20
\item  Regular slope of 1 m along the model (between 1 and 0 for the main channel)
\item  The bank level is 5 m upstream and the major bed grow regularly from 5 to 10 m
\end{itemize}

The Figure~\ref{fig:bridge:bathy} shows the mesh of the study.
\begin{figure}
\centering
\includegraphicsmaybe{[width=.8\textwidth]}{../img/Bathy.png}
\caption{Topography of the model}\label{fig:bridge:bathy}
\end{figure}

\section{Physical Parameters}

Turbulence: Model of constant viscosity with velocity diffusivity = 1 m${}^{2}$/s

\section{Numerical Parameters}

Type of advection:
\begin{itemize}
\item  edge-based N-scheme on velocities (scheme n$\mathrm{{}^\circ}$14)
\item  conservative + modified SUPG on depth (mandatory)
\item  edge-based N-scheme on tracer (scheme n$\mathrm{{}^\circ}$14)
\end{itemize}

Type of element:
\begin{itemize}
\item  Linear triangle P1 for h and velocities
\end{itemize}

\begin{itemize}
\item  Conjugate gradient solver
\item  Accuracy = ${10}^{-8}$
\end{itemize}

Tracer:

\begin{itemize}
\item  initial concentrations : no tracer in the domain
\item  Conjugate gradient solver (default value)
\item  Accuracy = ${10}^{-10}$
\end{itemize}

Time data:

\begin{itemize}
\item  Time step = 1 sec.
\item  Simulation duration = 45000 sec.
\end{itemize}

\section{ Results}

The Figure~\ref{fig:bridge:profile} illustrate the results on a longitudinal profile
 located in the middle of the main channel at 3 significant time-steps
\begin{figure}
\centering
\includegraphicsmaybe{[width=.8\textwidth]}{../img/FreeSurface_long_profile.png}
\caption{Free Surface profile along the channel}\label{fig:bridge:profile}
\end{figure}

During the first part of the simulation (before 6100 s), the flow stay in the minor bed 
as the culverts have sufficient capacities.
Velocities are regular except locally at the location of culvert extremities.
The tracer tends to an equilibrium in all the minor bed
See (Figure \ref{fig:bridge:surf0})

With the increase of the discharge, the culverts reach a saturation and then the level upstream
 increase until the overflow appears and a new equilibrium exists. This is illustrated
 by results at 17400 s on (Figure \ref{fig:bridge:surf1})
 
When the discharge decrease, there is a progressive diminution of the water lever upstream
 the dike. The simulation recover an equilibrium similar to the first one observed
 See (Figure \ref{fig:bridge:surf2})

\section{ Conclusions}

Bridge in TELEMAC-2D could be take in account by the combination of a local dike (which allow
 to represent the overflowing part) and some hydraulic structures which link the upstream and
 downstream of the dike

\section{ Figures}

\begin{figure}
\centering
 \includegraphicsmaybe{[width=1.\textwidth]}{../img/FreeSurface_0.png}
 \includegraphicsmaybe{[width=1.\textwidth]}{../img/Velo_0.png}
 \includegraphicsmaybe{[width=1.\textwidth]}{../img/Tracer_0.png}
 \caption{Map of the free surface elevation, the velocity and the tracer at t=6000s.}\label{fig:bridge:surf0}
\end{figure}

\begin{figure}
\centering
 \includegraphicsmaybe{[width=1.\textwidth]}{../img/FreeSurface_1.png}
 \includegraphicsmaybe{[width=1.\textwidth]}{../img/Velo_1.png}
 \includegraphicsmaybe{[width=1.\textwidth]}{../img/Tracer_1.png}
 \caption{Map of the free surface elevation, the velocity and the tracer at t=17400s.}\label{fig:bridge:surf1}
\end{figure}

\begin{figure}
\centering
 \includegraphicsmaybe{[width=1.\textwidth]}{../img/FreeSurface_2.png}
 \includegraphicsmaybe{[width=1.\textwidth]}{../img/Velo_2.png}
 \includegraphicsmaybe{[width=1.\textwidth]}{../img/Tracer_2.png}
 \caption{Map of the free surface elevation, the velocity and the tracer at t=45000s.}\label{fig:bridge:surf2}
\end{figure}

